\section*{Discussion}

Closed-loop optogenetic stimulation, driven by real-time wide-field imaging, held mesoscale cortical activity at a target more accurately and more repeatably than open-loop stimulation (Figure~\ref{fig:figure3}H,I), rejected part of a disturbance predicted from the opposite hemisphere (Figure~\ref{fig:figure4}G), and left an error that depended on brain state (Figure~\ref{fig:figure4}D). All of this rests on a controller whose structure and feedforward gain were set from the trial-averaged response to stimulation, with feedback gains tuned empirically, although single trials departed widely from that average.

\subsection*{A monotone trial-averaged response suffices to design a stabilizing controller}

We designed the controller for the trial average, and two of its properties carry the design. The fitted impulse-response models were stable, though their slowest mode is poorly constrained (Figure~\ref{fig:figure2}F), and evoked suppression grew monotonically, and approximately linearly, with laser power (Figure~\ref{fig:figure2}B). Monotonicity gives each feasible reference exactly one steady-state light drive that holds it, which is what integral action needs. For an exponentially stable plant whose steady-state input--output map is monotone, with slope bounded away from zero, a proportional--integral controller of matching sign and sufficiently low gains tracks constant references and rejects constant disturbances without a dynamic model of the plant~\citep{desoer1985tracking}. For integral control of a linear plant the conclusion survives a saturating, nondecreasing actuator and input delay~\citep{logemann2000absolute}, both features of our loop, although no single result covers proportional--integral control, saturation, delay and sampling together (Remark~\ref{rem:lowgain}). The integrator need not know the equilibrium input: its state stops changing only when the error is zero, so if the loop settles, it settles on the reference~\citep{francis1976internal}. Proportional feedback alone would leave part of any offset that the feedforward gain did not remove.

Variability between sessions therefore changes how high the gains can be set, not whether a stabilizing controller exists. With a single input and a stable plant, a stabilizing integral controller exists whenever the steady-state gain is nonzero, and one of fixed sign serves every plant whose gain has the matching sign~\citep{morari1985robust}. The variation we found between sessions, in gain (Figure~\ref{fig:figure2}B), response shape (Figure~\ref{fig:figure2}E) and slow rebound (Figure~\ref{fig:figure2}F), changes how fast the loop can settle, not the sign the controller needs, which is why we adjusted the gains at the start of every session. The guarantee is asymptotic and nominal (Remark~\ref{rem:lowgain}), so it justifies the controller's structure, not the error over 3\,s trials. Stability itself is empirical: no session showed sustained oscillation, the only instability a stable plant under a bounded command can show. The guarantee also treats the plant as fixed, and stability at every frozen state carries over only if the plant varies slowly enough~\citep{desoer1969slowly}, whereas brain state changes over seconds, not much slower than the loop settles. Finally, the sign of the map was measured, over our powers and in four sessions, not guaranteed, although interneuron drive suppresses excitatory activity in a graded way~\citep{li2019spatiotemporal}, and paradoxical responses of inhibition-stabilized cortex were absent on average with viral expression in parvalbumin interneurons~\citep{sanzeni2020inhibition}.

\subsection*{Linear on average, variable on every trial}

Linearity is a property of the average, not of single trials. A low-order linear model predicted held-out trial averages with a pooled median $R^2$ of 0.85 (Figure~\ref{fig:figure2}C,D), yet single trials scattered widely around the average, and a model fitted in one session described another only inconsistently (Figure~\ref{fig:figure2}E). At the macroscale, linear models describe aggregated brain activity at least as well as nonlinear ones, which their authors attribute to averaging and noise~\citep{Nozari2024}; trial averaging is one more averaging step of the same kind. We therefore treat the system from stimulation site to readout as nonlinear, driven by inputs from other regions that we do not observe, and dependent on a brain state that changes unmeasured, which makes it parameter dependent and non-stationary. Its trial average is the nominal plant we design on, and its single-trial departures are the disturbance that feedback must reject. Condition-averaged trajectories mix input-driven and recurrent contributions, which single-trial residuals can separate~\citep{galgali2023residual}.

Published models of cortical dynamics either fit single trials with inputs inferred after the fact or forecast them for barely a second, and both fall short for control for the same reason: the drive from other regions is unknown when a controller must act. Their fit over-approximates what is knowable in real time. Latent-dynamics models infer the inputs that perturb each trial with access to the whole trial~\citep{pandarinath2018lfads}, and data-constrained networks attribute a region's activity to currents from other recorded regions after the fact~\citep{perich2026curbd}. On spontaneous wide-field activity, the best of a dozen forecasters, each predicting a region from its own past, stays informative for about 1.5\,s~\citep{lu2025benchmarking}, and the authors of a whole-brain benchmark expect an irreducible error floor from unobserved processes and uncontrolled stimuli~\citep{lueckmann2025zapbench}. Leaving the inputs out of the fit does not help, because input-driven structure can then be mistaken for intrinsic dynamics~\citep{vahidi2024intrinsic}. Our own data show where the limit lies. The rest of cortex reveals much of this drive as it arrives: a stimulation-blind predictor from the contralateral hemisphere explained the spontaneous signal at the readout with a median held-out $R^2$ of 0.76 (Figure~\ref{fig:figure4}F). It reveals little of it in advance: a vector autoregressive model of the controlled site, 100 cortical sites and the laser predicted the controlled signal 200\,ms ahead with held-out $R^2 = 0.41$, against 0.37 for the site and laser alone, though the cortical sites added more at longer leads (Figure~\ref{fig:mpc_direct}H; Figure~\ref{fig:supp_varx}E). The drive is observable but not forecastable, which is the situation feedback is built for. This bounds only linear models of the observed cortex, but none of the forecasters above even takes the stimulation command that a controller requires.

For control, the model worth having is one that admits analysis. A low-order input--output model turns stability, admissible gain, delay-limited bandwidth and sensitivity to model error into computable quantities~\citep{astrom2008feedback,skogestad2005multivariable}, which a black-box forecaster, however accurate, does not; here we computed the feedforward gain from the dose--response relation and tuned the feedback gains by hand (Materials and Methods). Other neural systems have used input--output models of stimulation to design controllers~\citep{Yang2018} and to predict network responses~\citep{Yang2021}.We modelled only the map from laser to readout and left the rest to feedback; the uncertainty this leaves is acceptable as long as it does not destabilize the loop.

State dependence belongs to the unstimulated cortex, but it appears in stimulus responses, which is why state-resolved responses are worth characterizing. Single-trial deviations from the trial average fell in the highest-motion quartile and rose with absolute 2--4\,Hz power, and ongoing activity in a stimulus-free window did the same at a similar magnitude (Figure~\ref{fig:figure2}G); relative to that window, response variability barely changed across quartiles (Figure~\ref{fig:state_quartile}). The state dependence we see is therefore carried mainly by the activity the stimulus lands on, as single-trial responses ride on the ongoing activity that precedes them~\citep{Arieli1996,Shimaoka2019}, although a change in the stimulus effect itself is not excluded~\citep{Mateo2011}. The distinction matters for future models: state that adds an offset is a disturbance that feedback subtracts, whereas state that changes gain~\citep{Goris2014} is a model error that feedback must tolerate. Switching linear dynamical systems~\citep{linderman2017rslds}, extended with a stimulation input and fitted to state-resolved trials such as ours, would keep that analysability regime by regime.

\subsection*{Feedback partly cancels activity shared with the opposite hemisphere}

The controller cancels part of the input the controlled site receives from the rest of cortex, as far as the opposite hemisphere reveals it, and holds the site near the target. Measured against the contralaterally predicted disturbance, the open loop left more energy about the reference and the closed loop less (median ratios 1.25 and 0.72 in the same eleven sessions; Figure~\ref{fig:figure4}G, Table~\ref{tab:stat_fig4g}). Most of the difference was the standing offset, which a fixed input leaves through its calibration rather than as part of the contralateral drive: it fell from 35\% of the residual energy to 7\%, whereas the fluctuation fell only from 0.66 to 0.54 of the disturbance. This is what a proportional--integral loop should do. Its integrator is an internal model of constant signals, so it cancels a constant error, and the loop attenuates other disturbances only at frequencies where its loop gain, the amplification once around the loop, is large~\citep{francis1976internal,astrom2008feedback}. Interhemispheric drive is a natural target for such cancellation: in premotor cortex, activity disrupted by unilateral silencing is restored by input from the other hemisphere~\citep{li2016robust}, so an open-loop perturbation of one side is likely opposed by the other. The contralateral split is, however, a linear projection rather than a causal decomposition. Optogenetically evoked interhemispheric connectivity is weaker than resting-state correlations where those correlations are strong~\citep{bauer2018effective}, so the shared component may partly reflect common subcortical, neuromodulatory or movement-related drive.

Because the loop holds a site near a set level rather than fixing the light dose, it constrains the perturbation on each trial far more tightly than a fixed dose does (across-trial variance 46\% lower from 1 to 3\,s; Figure~\ref{fig:figure3}I), which is what causal tests of downstream effects need. Open-loop manipulations push the network into unusual states~\citep{Wolff2018}, evoke a different response on every trial (Figure~\ref{fig:figure2}G) and act off-target downstream~\citep{otchy2015acute}. The logic is that of the voltage clamp, in which the current needed to hold the membrane at a command potential measures the ionic current the membrane passes there~\citep{hodgkin1952measurement}, carried to networks by the optoclamp~\citep{newman2015optogenetic}. The laser command needed to hold the readout should likewise track the drive the loop opposes, a relation we did not test; if it does, the command is itself a readout of the cancelled input. With a site held at a chosen level, its influence can be studied as a function of a controlled upstream value, across the rest of dorsal cortex in the field of view, in striatum~\citep{Peters2021} or in behaviour~\citep{ZatkaHaas2021}. Unlike the voltage clamp, there is no space clamp: surrounding and contralateral tissue are not held, and the regulated quantity is the readout kernel, not the region.

\subsection*{Reachability and predictability bound what feedback can do}

The actuator sets which references the loop can reach. With one light spot that can only suppress, a reference is reachable only if its steady-state suppression lies within what the allowed commands can deliver, the feasibility condition of Remark~\ref{rem:lowgain}. This is a steady-state condition rather than controllability, although a single input of one sign also cannot steer a linear system with a real mode in every direction~\citep{brammer1972controllability}. Our data are consistent with this bound. Commands reached the ceiling in seven of the 13 sessions recorded in 2025 (Materials and Methods), the closed-loop average never reached the reference in the two sessions excluded from the disturbance analysis for a standing offset, and about half of the squared error left by the model-predictive replay with a sinusoidal reference fell in the half-cycles where the reference rose, when the laser was mostly off. Slow integral action within a trial, or the clamp on the integral term, could also explain each shortfall. A second, excitatory actuator would make both directions reachable.

Within the reachable range, the slopes of Figure~\ref{fig:figure4}D separate disturbances by whether the loop has gain against them. The open-loop slope, predictability, asks whether a state enlarges the disturbance, and the closed-loop slope, regularizability, whether the loop still holds the output. The ratio of closed- to open-loop error is a broadband analogue of the sensitivity function, the fraction of a disturbance that a loop leaves, averaged over frequency and weighted by the open-loop error spectrum~\citep{astrom2008feedback}; attenuation, the change of this ratio with state (Table~\ref{tab:fig4_state_stats}), asks whether that fraction depends on state. Movement-related drive enlarged the disturbance but was rejected: open-loop error rose with motion, closed-loop error did not, and the fraction left by the loop fell by 9\% per SD of motion. Slow, large-amplitude fluctuations were reduced but not held: absolute 1--4\,Hz power raised closed-loop error as much as open-loop error, and the fraction left rose by 8\% per SD. The initial deviation behaved as a linear loop predicts: it dominated closed-loop error in the first second (Figure~\ref{fig:figure4}C) but left the ratio from 1 to 3\,s unchanged, so the loop removed the same fraction of a large deviation as of a small one.

Loop timing, together with the lag of the cortical response, is one limit on slow fluctuations. An end-to-end delay of at most 47\,ms, about half that on average (Figure~\ref{fig:figure1}F), plus the 14\,ms average wait for the next 35\,Hz frame, gives an effective delay of 37--61\,ms, which by the crossover criterion $\omega_c \tau \lesssim 1$\,rad caps the bandwidth of any controller at about 2.6--4\,Hz~\citep{skogestad2005multivariable}; stricter rules of thumb place it lower~\citep{astrom2000limitations}. Above crossover, where loop gain falls below one, feedback attenuates little, and near crossover it amplifies, the price that Bode's sensitivity integral exacts for attenuation at lower frequencies~\citep{astrom2008feedback}. The cap counts only the electronics; the cortical response adds about 60 degrees of lag at 1\,Hz (Figure~\ref{fig:figure5}), so the loop that ran probably crossed over lower, perhaps near 1--2\,Hz, which would put most of the 1--4\,Hz band at or above crossover, though we did not measure it. Two other accounts remain open. Large slow excursions drive the command into its one-sided ceiling, where no linear argument applies, and band power was measured on the controlled signal during stimulation, so the loop may partly shape its own closed-loop slope. Movement and slow power are also anticorrelated in awake cortex~\citep{McGinley2015waking}, so the two state effects need not be independent.

\subsection*{Prediction removes error that feedback leaves}

What feedback leaves is bounded by what can be predicted. For a linear plant with dead time and an additive stationary disturbance, no causal controller can push the output variance below that of the part of the disturbance that cannot be predicted over the dead time~\citep{astrom1970introduction}. State dependence, which makes our plant time-varying and the disturbance non-stationary, turns this floor into a benchmark, but the conclusion stands: the error a proportional--integral loop leaves generally contains a predictable part, and removing it takes a controller built on a model of the disturbance, which is what prediction supplies.

Predicting the reference is the simplest case. With a 1\,Hz reference, preview cancelled most of the cortical phase lag but barely changed RMSE, because the constant offset remained the largest error term (Figure~\ref{fig:figure5}), so studies with time-varying targets should report offset, drive-frequency and broadband errors separately. A fixed five-sample lookahead is not optimal preview~\citep{tomizuka1975discrete}, and under feedback, which had already removed part of the lag, ours overcorrected.

Predicting the disturbance is the larger gain. In a replay of recorded trials, a model-predictive controller built on the vector autoregressive model lowered the median per-trial error over 1--3\,s from 1.46 to $0.88\,\%\,\Delta F/F$, on 107 of 108 trials, and across-trial variance by 37\% (Figure~\ref{fig:mpc_direct}C--F). Given the future residuals it reached 0.27, so most of what it left was activity it could not predict. Model-predictive control adds three things to proportional--integral feedback: it anticipates the predictable part of the disturbance, treats the one-sided, bounded laser as an explicit constraint rather than a saturation applied afterwards~\citep{mayne2000constrained}, and previews the reference in its cost~\citep{rawlings2017mpc}. With the sinusoidal reference it kept the phase lag near zero without a tuned lookahead (Figure~\ref{fig:mpc_direct}M), but the gain was smaller (3.00 to $2.67\,\%\,\Delta F/F$ over 4\,s), and even perfect prediction reached only 0.77 times the error of feedback with preview. The replayed cortex shares the controller's model class and responds to the laser identically in every state, so these values bound rather than measure what a live controller would achieve.

A live predictive controller should keep integral action. One built on an imperfect model generally leaves a steady offset unless its model carries integrating disturbances estimated by an observer, and even then only if the constraints leave the target reachable~\citep{pannocchia2003disturbance}. Offset was the largest error term with the sinusoidal reference (Figure~\ref{fig:figure5}) and a third of the open-loop residual with the constant one (Table~\ref{tab:stat_fig4g}), so offset-free predictive control, or a predictive outer loop that sets references for the present regulator, is the natural design. Predictive control of neural activity has so far been shown mainly in simulation~\citep{chang2020seizure,steffen2024mpcdbs}, and model-based optimal feedback runs live at the single-neuron scale~\citep{bolus2021state}; our replay, on real mesoscale data, sits between the two. Beyond anticipating the disturbance and respecting the one-sided laser, a live predictive controller could use brain state rather than average it away, for example through linear predictors in a lifted state space, learned from each session's data, that capture nonlinear dynamics while keeping the control problem convex~\citep{korda2018koopman}.

\subsection*{Limitations}

The controller fed back on the blue channel without hemodynamic correction, which needs the interleaved violet frame and cannot be computed within one control interval. Every closed-loop analysis, including the band-power state measures and the contralateral predictor, therefore describes the signal the loop regulated rather than purely neural activity~\citep{valley2020separation}. The impulse analysis, on the corrected signal, showed the absolute 2--4\,Hz effect in every session, but we cannot apportion the closed-loop effect between neural and vascular sources. Inhibition spreads beyond the stimulation spot (Figure~\ref{fig:spatial_spread}A,B;~\citealp{li2019spatiotemporal}), and the readout lies next to it, so the loop partly regulates that spread. The sample is small and the scope narrow: four impulse sessions from three mice, 15 sessions from four mice for the open- versus closed-loop comparison, eleven from three mice for disturbance rejection, three sinusoidal sessions from one mouse and one session for each replay, all with inhibitory stimulation of one somatosensory site. Interleaving open- and closed-loop trials guards against drift in brain state between conditions but cannot remove it, and whether these results carry over to excitation, to other areas or to larger perturbations is untested.
